\section{Precision Ladder}
\label{app:precision-ladder}

\subsection{Implemented routes}

The copied scheduler contains a ladder of modes so format changes can be
isolated without changing the BF16 control.  Table~\ref{tab:precision-ladder}
states the mathematical intent rather than every historical compile-time
probe.

\begin{table}[htbp]
\centering
\caption{Implemented backward precision ladder.}
\label{tab:precision-ladder}
\scriptsize
\begin{tabularx}{\textwidth}{c l l l X}
\toprule
Mode & Q/K score & $dP,dV$ operands & $dS,dQ,dK$ & Status \\
\midrule
0 & BF16 & BF16 & BF16 & copied V382 control \\
1 & E4M3 & BF16 & E4M3 dS, E4M3 $dQ/dK$ & native-FP8 diagnostic \\
2 & E2M1 & BF16 & E4M3 dS, mixed E4M3$\times$E2M1 & first FP4 rung \\
5 & adaptive E2M1 & E4M3 $dP,dV$ & E4M3 dS, mixed dQ/dK & retained winner \\
12 & E2M1 & mixed E4M3/MXFP4 dP; E4M3 dV &
E4M3 dS, mixed dQ/dK & fixed-scale speed frontier \\
12+A & adaptive E2M1 & mixed E4M3/MXFP4 dP; E4M3 dV &
E4M3 dS, mixed dQ/dK & lifetime-safe H24 frontier \\
11 & E2M1 NVFP4 & MXFP4 dP and dV &
MXFP4 dS and scaled MXFP4 dQ/dK & experimental pure route \\
\bottomrule
\end{tabularx}
\end{table}

\subsection{Fixed mixed-dP equation}

The faster mode-12 route keeps the first K64 of V and $G$ in E4M3 and packs
the second K64 as fixed-scale E2M1.  Abstractly,
\begin{equation}
  D_P^{\mathrm{mixed}}
  =
  \mma_{\mathrm{E4M3}}(\bar G_{0:64},\bar V_{0:64}^{\mathsf T})
  +\gamma\,
  \mma_{\mathrm{E2M1}}(\bar G_{64:128},\bar V_{64:128}^{\mathsf T}),
  \label{eq:mixed-dp}
\end{equation}
where $\gamma$ is folded into a reused fixed scale page.  dV remains the
all-E4M3 operation in Equation~\ref{eq:represented-dv}.  A reusable V row
contains 64 E4M3 bytes, 32 packed E2M1 bytes, and 32 padding bytes.  Forward
can emit this row once so backward does not repay the V conversion.

This route is 3.3--5.3\% faster than the retained winner on the broad matrix,
but its calibrated dQ/dK cosine is approximately 0.983 rather than 0.991.
The scale is fixed and distribution-dependent.

Mode 12+A supplies the robust per-head Q/K record to the same mixed-dP
operator.  Its score and dQ/dK scale page is distinct from the fixed dP tail
word, resolving the lifetime bug in the original attempt.  It is 2.3--2.7\%
faster than adaptive all-E4M3 dP/dV at S4096/H24 but 1.3--1.7\% slower at
S8192/H8.  Because mixed dP still controls its component error, calibrated
dQ/dK cosine remains approximately 0.983.

\subsection{MXFP4 block representation}

For a 32-value block B with magnitude $a_B$, choose a power-of-two scale
\begin{equation}
  s_B=2^{e_B},\qquad
  e_B\in
  \left\{
    \left\lfloor\log_2(a_B/6)\right\rfloor,
    \left\lceil\log_2(a_B/6)\right\rceil
  \right\},
  \label{eq:mxfp4-block-scale}
\end{equation}
subject to E8M0 range and the selected rounding policy.  The represented
block is
\begin{equation}
  \bar X_B=Q_{\mathrm{E2M1}}(X_B/s_B),\qquad
  \widehat X_B=s_B\bar X_B.
  \label{eq:mxfp4-block}
\end{equation}
The tensor instruction receives swizzled 32$\times$16 scale pages in TMEM.

\subsection{Experimental pure-MXFP4 backward}

The pure route applies Equation~\ref{eq:mxfp4-block} to V, $G$, P, and dS.
At the represented-operator level:
\begin{align}
  D_P^{(4)}
    &=\mma_{\mathrm{MX4}}(\bar G_4,\bar V_4^\mathsf T;
       s_G,s_V),\\
  \widetilde{dV}^{(4)}
    &=\mma_{\mathrm{MX4}}(\bar P_4^\mathsf T,\bar G_4;
       s_P,s_G).
  \label{eq:pure-dp-dv}
\end{align}
The raw MXFP4 dP accumulator carries the implementation's standard
width-six reconstruction factor:
\begin{equation}
  \widetilde{dP}^{(4)}=\frac{1}{36}D_P^{(4)}.
  \label{eq:mxfp4-dp-correction}
\end{equation}
The producer folds $1/36$ into centering instead of materializing a corrected
dP tile:
\begin{equation}
  dS^\star
  =\widehat P_4\odot
   \left(\frac{1}{36}D_P^{(4)}-r\mathbf1^\mathsf T\right).
  \label{eq:pure-ds-prequant}
\end{equation}
It then emits two independently scaled MXFP4 orientations,
\begin{align}
  (\bar dS_K,s_{dS,K})&=Q_{\mathrm{MX4}}(dS^\star)
    &&\text{for }dK,\\
  (\bar dS_Q,s_{dS,Q})&=Q_{\mathrm{MX4}}((dS^\star)^\mathsf T)
    &&\text{for }dQ,
  \label{eq:pure-dual-ds}
\end{align}
and consumes compact Q/K operands:
\begin{align}
  \widetilde{dQ}^{(4)}
  &=\alpha\,
    \mma_{\mathrm{MX4}}(\bar dS_Q,\bar K_4;
      s_{dS,Q},s_K),\\
  \widetilde{dK}^{(4)}
  &=\alpha\,
    \mma_{\mathrm{MX4}}(\bar dS_K,\bar Q_4;
      s_{dS,K},s_Q).
  \label{eq:pure-dq-dk}
\end{align}

P uses a forward-style log-domain producer.  For one probability block,
\begin{align}
  x_{ij} &=
    (\widetilde Z_{ij}-L_i)\log_2e-e_B+\log_2 6,\\
  \bar P_{4,ij} &=Q_{\mathrm{E2M1}}(2^{x_{ij}}),\qquad
  \widehat P_{4,ij}=\frac{s_B}{6}\bar P_{4,ij}.
  \label{eq:pure-log-p}
\end{align}
The block maximum that selects $e_B$ is shared with quantization, avoiding a
second full probability pass.

The initial 148-tile ceiling validated every active matrix output bit-exactly
against products of the captured represented operands.  Its recommended
NV-dP/MX-dS stage-6 point measured 130.142 $\mu$s for 148 independent tiles.
The completed long-sequence mode 11 now also consumes projection-native Q/K/V
and producer-native dO payloads.  Table~\ref{tab:pure-mxfp4-broad} reports the
prepared-kernel comparison from one fixed-seed calibrated sweep.

\begin{table}[htbp]
\centering
\caption{Prepared backward precision frontier.  Times are median milliseconds;
the final column is aggregate dQ/dK/dV cosine of pure MXFP4 versus BF16.}
\label{tab:pure-mxfp4-broad}
\scriptsize
\begin{tabular}{l rrrrr r}
\toprule
Shape & BF16 & FP4+FP8 & mixed dP & adaptive FP4+FP8 & pure MXFP4 & pure cosine \\
\midrule
S4096/H24  & 0.707936 & 0.647104 & 0.613664 & 0.660096 & 0.637120 & 0.984797 \\
S4096/H64  & 1.655648 & 1.513952 & 1.432832 & 1.539360 & 1.482304 & 0.984614 \\
S8192/H8   & 0.869984 & 0.758656 & 0.745664 & 0.773472 & 0.781920 & 0.984653 \\
S8192/H24  & 2.212064 & 1.923200 & 1.885216 & 1.956544 & 1.971136 & 0.985197 \\
S8192/H64  & 5.558176 & 4.871296 & 4.739136 & 4.940192 & 4.958176 & 0.985087 \\
S16384/H24 & 7.758656 & 6.510784 & 6.536256 & 6.610144 & 6.874848 & 0.985217 \\
S16384/H64 & 20.275393 & 17.163391 & 17.102079 & 17.419872 & 17.997663 & 0.985224 \\
\bottomrule
\end{tabular}
\end{table}

Pure MXFP4 is 10.0--11.4\% faster than BF16 in this matrix and beats the
ordinary fixed FP4+FP8 route on the two S4096 rows.  It does not beat the
mixed-dP route on any row, and it loses to adaptive FP4+FP8 once sequence is
8192 or longer.  More importantly, its dQ/dK component cosine is only about
0.864--0.887, while aggregate cosine near 0.985 is dominated by the more
accurate dV component.  The stable long-sequence accuracy-oriented API
therefore remains mode 5; mode 11 is a validated aggressive endpoint.

\subsection{Why pure FP4 is not the headline}

The pure route eliminates more operand bytes but adds block maxima, E8M0 scale
publication, two dS scale orientations, and additional TMEM scale lifetimes.
Backward's scalar and synchronization work is already exposed, so replacing
E4M3 instructions with MXFP4 does not automatically translate the peak
tensor-core ratio into wall time.  Fusing Q/K quantization into projection
removes 69.6--97.2~$\mu$s from two composed checks, but it does not shorten
those in-kernel maxima and scale handoffs.  Its numerical contract is also
more aggressive than the retained approximately 0.991 dQ/dK point.  It
remains the right research ladder, but not the current default.

Replacing either scale family with NVFP4 K16/E4M3 pages does not change that
conclusion.  NVFP4 dS with MXFP4 Q/K improves dQ/dK cosine to approximately
0.969 but becomes slower than BF16; NVFP4 Q/K loses speed and accuracy; and
the combined NV/NV diagnostic is non-finite.  The retained pure specialization
therefore keeps MXFP4 dS and MXFP4 Q/K scale pages.
